% Generated draft: requires mathematical and source review.
\documentclass[11pt]{article}
\usepackage[margin=25mm]{geometry}
\usepackage{fontspec}
\setmainfont{Latin Modern Roman}
\usepackage{amsmath,amssymb,mathtools,unicode-math}
\setmathfont{Latin Modern Math}
\usepackage{enumitem,needspace,xurl,hyperref}
\setlength{\parindent}{0pt}
\setlength{\parskip}{5pt}
\setlength{\emergencystretch}{3em}
\newcommand{\N}{\mathbb{N}}
\newcommand{\Z}{\mathbb{Z}}
\newcommand{\Q}{\mathbb{Q}}
\newcommand{\R}{\mathbb{R}}
\newcommand{\C}{\mathbb{C}}
\newcommand{\F}{\mathbb{F}}
\newcommand{\A}{\mathbb{A}}
\newcommand{\PP}{\mathbb{P}}
\newcommand{\E}{\mathbb{E}}
\newcommand{\eps}{\varepsilon}
\newcommand{\dd}{\,\mathrm{d}}
\begin{document}
\section*{UnsolvedMath ID: 30001409}
\textbf{Source problem code:} catalogue entry
\subsection{Definitions and mathematical statement}
The supplied record presents three related questions about affine surface area, but does not specify enough hypotheses to determine a unique mathematical formulation. The following notation describes the questions while keeping those omissions explicit.

For an integer \(n\geq 2\), write \(B^n=\{x\in\R^n:\|x\|_2\leq 1\}\). A convex body is a compact convex subset of \(\R^n\) with nonempty interior. For such a body \(K\) and \(\varepsilon\geq 0\), its outer parallel body is
\[
K+\varepsilon B^n=\{x+\varepsilon y:x\in K,\ y\in B^n\}.
\]
These are conventional meanings of the notation; the source does not explicitly give the permitted dimensions, the class of bodies, or the parameter range.

For orientation, if \(K\) has a \(C^2\) boundary with positive Gaussian curvature \(\kappa_K\), the usual affine surface area is
\[
\Omega(K)=\int_{\partial K}\kappa_K(x)^{1/(n+1)}\,\mathrm{d}\mathcal H^{n-1}(x),
\]
where \(\mathcal H^{n-1}\) is Euclidean surface measure. This formula describes the classical smooth setting only. The supplied record does not specify the definition on nonsmooth bodies or state that its questions are restricted to this smooth setting.

\begin{enumerate}
\item Determine whether affine surface area admits a Steiner formula for \(K+\varepsilon B^n\), and determine the corresponding assertion for its ``maximal variant.'' The source does not define this maximal variant. It also does not specify whether a Steiner formula means a finite polynomial identity, a convergent series, an asymptotic expansion, or another prescribed representation, nor does it specify its coefficients or range of validity. Thus the target identity remains unresolved at the level of formulation.
\item Determine explicitly the functional called a valuation in the source,
\[
\Phi(K)=\int\Omega(K|E)\,\mathrm{d}E.
\]
Here \(K|E\) conventionally denotes orthogonal projection onto a linear subspace \(E\), with affine surface area then computed in \(E\). However, the source does not specify whether this projection convention is intended, the dimension of \(E\), the integration space, the measure or its normalization, or the applicable definition of affine surface area in the projected dimension. More explicitly, the displayed expression is only the schematic formula
\[
\Phi(K)=\int_{\mathcal E}\Omega(K|E)\,\mathrm{d}\mu(E),
\]
with \(\mathcal E\) and \(\mu\) unspecified. Calling a functional a valuation conventionally means that
\[
\Phi(K\cup L)+\Phi(K\cap L)=\Phi(K)+\Phi(L)
\]
whenever all four sets belong to its domain. The domain and any convention for lower-dimensional or empty intersections are not supplied here; the valuation designation is retained as a source description.
\item Decide whether affine kinematic formulae exist for the affine-surface-area setting under discussion. The source does not specify the transformation group, averaging measure, operation on bodies, integrand, or desired right-hand side. Consequently this is a research question with an unspecified target family of identities, rather than a uniquely quantified existence statement.
\end{enumerate}

The conventional definitions above do not impose missing hypotheses on the original questions. A complete precise statement requires recovery of the intended definitions and quantifiers from the cited report.
\subsection{Short English statement}
The record asks for Steiner formulae for affine surface area and its maximal variant, an explicit evaluation of an averaged projection functional, and affine kinematic formulae. The maximal variant, averaging data, admissible bodies, and intended identities are unspecified. The supplied status is \emph{partially solved}; this draft preserves that label without independently verifying it.
\subsection{Sources}
\begin{itemize}
\item Local UnsolvedMath registry: supplied complete source record, permanent ID \(30001409\), title \emph{Steiner and Kinematic Formulas for Affine Surface Area}, exact source problem code \texttt{catalogue entry}, in the Oberwolfach Reports Open Problems collection. The record gives the proposed year as 2010 and names no individual proposer.
\item Supplied original reference: \emph{Mini-Workshop: Valuations and Integral Geometry} (2010), Oberwolfach Reports, pp.~141--178, DOI \texttt{10.4171/owr/2010/04}. \url{https://doi.org/10.4171/owr/2010/04}. The supplied background describes these reports as documenting mathematical workshops and their open-problem discussions.
\item Supplied literature source, labelled ``later additive kinematic formulas'': \url{https://www.cambridge.org/core/journals/canadian-journal-of-mathematics/article/additive-kinematic-formulas-for-convex-functions/4B02C40FE7A094A91711AFBBF4DD77FF}.
\item Supplied literature assessment dated exactly \texttt{2026-08-21}: status \emph{partially solved}. The supplied assessment states that the report records partial answers for the affine-surface-area Steiner question and asks for broader kinematic formulae, and that subsequent work supplies related additive kinematic formulae but not the full collection stated here. This is supplied dated curation derived from the supplied OWR curation, not an independently verified literature assessment; its time-sensitive status and cited claims require independent verification.
\item The registry labels the statement \texttt{corrected\_verified}, but its accompanying verification says only that TeX notation was enclosed in math delimiters without changing mathematical content. This does not establish mathematical correctness or resolve the missing formulation data. Neither linked source was visited for this draft.
\end{itemize}
\end{document}
